\section{元探索：学习如何探索}

\subsection{什么是元探索}

\begin{importantbox}{Meta-Exploration}
元探索（meta-exploration）将元学习的思想应用于探索问题：不是手工设计探索策略，而是\textbf{学习一个探索策略}。

在 meta-RL 的框架下：
\begin{itemize}
    \item 外层（meta）学习如何在新任务上有效探索
    \item 内层（task）利用探索收集的数据快速适应新任务
\end{itemize}
\end{importantbox}

\subsection{Meta-RL 中的探索挑战}

在 meta-RL 中，智能体在每个新任务上的前几次交互至关重要——它们决定了智能体能否快速推断出任务的特性。因此，元探索关注的是：\textbf{在有限的交互次数内，如何收集最有信息量的数据}。

\begin{knowledgebox}{与经典探索的区别}
经典探索关注单个任务中的 exploration-exploitation 权衡。元探索关注的是跨任务的探索策略学习——训练一个能在新任务上有效探索的策略。这是一个更高层次的优化问题。
\end{knowledgebox}

\begin{figure}[H]
\centering
\includegraphics[width=0.9\textwidth]{slides-images/slide-15.jpg}
\caption{当测试时仍需探索：厨房找食材、LLM 推理选策略都属于 meta-RL 式探索}
\end{figure}

\subsection{耦合问题：为什么端到端学探索经常失败}

如果探索和执行都只通过最终任务奖励来训练，就会出现典型的 chicken-and-egg 困境：探索没做好，执行阶段拿不到好 reward；执行阶段一直失败，探索阶段也就收不到足够清晰的训练信号。

\begin{figure}[H]
\centering
\includegraphics[width=0.9\textwidth]{slides-images/slide-16.jpg}
\caption{端到端 meta-learning with exploration 的 coupling problem}
\end{figure}

\subsection{DREAM 和后验采样方法}

讲者介绍了如何在 meta-RL 中实现有效探索：
\begin{itemize}
    \item 维护对任务身份的信念（belief）
    \item 使用 Thompson Sampling 的思想指导探索
    \item 外层 RL 训练使内层探索策略最大化信息增益
\end{itemize}

\subsection{替代策略一：后验采样、内在奖励与任务预测}

课程先回顾了几类已有思路：PEARL 用 posterior sampling，MAME 用 intrinsic rewards，MetaCURE 用 task dynamics/reward prediction。它们的共同点是：\textbf{不再把探索完全交给 end-to-end reward 信号，而是显式引入任务推断结构。}

\begin{figure}[H]
\centering
\includegraphics[width=0.9\textwidth]{slides-images/slide-17.jpg}
\caption{Alternative strategies：posterior sampling、intrinsic rewards、task dynamics prediction}
\end{figure}

\subsection{DREAM：把探索和执行显式解耦}

这节课真正的主角是 DREAM。它的关键想法不是直接学一个同时负责探索和执行的黑盒策略，而是分两步：
\begin{enumerate}
    \item 学会执行，同时找出真正决定任务差异的信息瓶颈表示。
    \item 再学一个探索策略，让它专门去恢复这些任务相关信息。
\end{enumerate}

\begin{figure}[H]
\centering
\includegraphics[width=0.9\textwidth]{slides-images/slide-20.jpg}
\caption{DREAM 的解耦思想：先识别任务相关信息，再训练探索去恢复这些信息}
\end{figure}

\begin{knowledgebox}{为什么信息瓶颈是必要的}
如果不压缩任务表示，探索策略很容易被与任务无关的装饰性信息带偏。DREAM 借助 bottlenecked representation，只保留 wall color、食材位置等真正影响执行策略的变量，让探索策略的奖励更聚焦于 ``找对信息''。
\end{knowledgebox}

\begin{figure}[H]
\centering
\includegraphics[width=0.9\textwidth]{slides-images/slide-22.jpg}
\caption{DREAM 的训练方式：执行策略学会依赖 bottlenecked task representation，探索策略学会恢复它}
\end{figure}

\subsection{把信息恢复直接变成探索奖励}

课程进一步把探索奖励写得很清楚：让探索策略收集的数据 $D_{\text{train}}$ 尽可能有助于预测任务表示 $z_i$。这样 reward 不再依赖最终任务是否成功，而依赖 ``本轮探索让我们对当前任务多了解了多少''。

\begin{figure}[H]
\centering
\includegraphics[width=0.9\textwidth]{slides-images/slide-23.jpg}
\caption{DREAM 的核心奖励：per-step information gain，而不是最终任务 reward}
\end{figure}

\subsection{信息瓶颈为什么值得单独建模}

如果任务表示里塞进了太多与执行无关的细节，探索策略就会去恢复错误的信息，例如装饰、纹理或无关背景。课程专门回顾 variational information bottleneck 的原因就在这里：通过噪声和正则约束，把表示压到只剩执行真正需要的那部分任务差异。

\begin{figure}[H]
\centering
\includegraphics[width=0.9\textwidth]{slides-images/slide-21.jpg}
\caption{Variational information bottleneck：压缩任务表示，只保留与执行有关的信息}
\end{figure}

\subsection{理论与实验结果}

slides 给出的结论很有代表性：在 bandit-like setting 中，DREAM 的样本复杂度优于 RL$^2$；在 3D visual navigation 这种稀疏奖励任务里，端到端方法会因为 coupling 问题显著掉队，而 DREAM 靠近似最优探索达到更高回报。

\begin{figure}[H]
\centering
\includegraphics[width=0.9\textwidth]{slides-images/slide-24.jpg}
\caption{理论分析：DREAM 在某些设置下保留最优探索性质，同时显著改进样本复杂度}
\end{figure}

\begin{figure}[H]
\centering
\includegraphics[width=0.9\textwidth]{slides-images/slide-25.jpg}
\caption{稀疏奖励 3D 视觉导航任务：先读标志，再绕障碍到达正确目标}
\end{figure}

\begin{figure}[H]
\centering
\includegraphics[width=0.9\textwidth]{slides-images/slide-26.jpg}
\caption{稀疏奖励 3D 视觉导航中的实验结果：DREAM 接近最优，端到端方法受 coupling 影响明显}
\end{figure}

\begin{figure}[H]
\centering
\includegraphics[width=0.9\textwidth]{slides-images/slide-27.jpg}
\caption{三类方案的总结：端到端、替代策略与解耦探索执行}
\end{figure}

\subsection{本章小结}

元探索将探索策略本身作为学习对象，通过在多个任务上训练来学习高效的探索行为。

